\documentclass[11pt]{article}
\usepackage{fontspec}
\usepackage[french]{babel}
\usepackage{amsmath}
\usepackage{amssymb}
\usepackage{amsthm}
\usepackage{mathrsfs}
\usepackage{tikz-cd}
\usepackage{enumitem}

\newcommand{\printedpage}[1]{\marginpar{\footnotesize p.~#1}}
\newcommand{\msp}[1]{\marginpar{\footnotesize [#1]}}
\newcommand{\Hom}{\operatorname{Hom}}
\newcommand{\Ker}{\operatorname{Ker}}
\newcommand{\Ob}{\operatorname{Ob}}

\begin{document}

\printedpage{19}
% head: A. Grothendieck : Dérivateurs, Ch. II, Édition des Universités de Montpellier II et Paris VII
% foot: Version 8 juillet 2001

On est ramené à prouver que cette dernière condition implique la condition d) de 2.6, et on peut prendre pour $S$ l'ensemble ordonné des ouverts connexes non vides de $S$, cela devient alors évident.

Ceci dit, il faut se rappeler que la théorie cohomologique des espaces topologiques irréductibles n'est nullement triviale - comme nous le montre FAC de \textsc{Serre}\,! D'autre part, c'est \textsc{Serre} qui m'avait fait observer que pour des coefficients \textit{constants,} la cohomologie d'un espace irréductible est triviale.

Dans le cas actuel, prouvons déjà, si $\mathcal{X}$ est 0-connexe et localement totalement 0-connexe, que tout revêtement de l'objet final de $\mathcal{X}$ est trivial. Il suffit de voir que si $U$ est 0-connexe, $U \longrightarrow e$ est un isomorphisme. Mais considérons $U \hookrightarrow U \times U$, on sait que c'est un sommande direct, et comme $U \times U$ est connexe, c'est un isomorphisme. Donc $U \longrightarrow e$ \texttt{[est un]} monomorphisme, donc sommande direct, et comme $U \neq \emptyset$, on a $U \longrightarrow e$ isomorphisme.

Reste à prouver que pour tout faisceau \msp{page 31} abélien constant $F$ sur $\mathcal{X}$, on a $H^i(X,F)=0$ pour $i>0$. On procède par récurrence sur $i$, c'est déjà fait pour $i=1$. Supposons que ce soit prouvé pour les $j<i$ ($i$ fixé) pour tout topos \textit{totalement} 0\textit{-connexe.} Soit $\xi \in H^i(\mathcal{X},F)$, soit $U \longrightarrow e$ \texttt{[un]} épimorphisme tel que $\xi|U=0$, et considérons la suite spectrale % [sic] H^i(X,F) -> H^i(\mathcal{X},F), plain X printed where the topos \mathcal{X} is meant elsewhere
\[
E_2^{pq} = H^p(C(U^*,\mathcal{H}^q(F))) \Longrightarrow H^*(\mathcal{X},F).
\]
On a par hypothèse totalement 0-connexe
\[
\pi_0(U^q) \simeq \pi_0(U)^q, \tag*{$(*)$}
\]
et les composantes connexes de $U^q$ sont totalement 0-asphériques,
\begin{center}
\begin{tikzpicture}[scale=0.8]
\draw[->] (0,0) -- (0,5);
\node[left] at (0,5) {$q$};
\draw[->] (0,0) -- (5,0);
\node at (4.3,-0.6) {$p$};
\node[left] at (0,4) {$i$};
\foreach \x in {0,1,2,3,4} {\fill (\x,4) circle (2pt); \fill (\x,0) circle (2pt);}
\end{tikzpicture}
\end{center}
donc par hypothèse de récurrence
\[
E_1^{pq} = 0 \qquad \text{si } 0<q<i,
\]
donc on trouve une suite exacte
\[
0 \longrightarrow H^i(C(U^*,F)) \longrightarrow H^i(\mathcal{X},F) \longrightarrow H^0(\ \underbrace{H^i(U^*,F)}_{\substack{\text{complexe}\\ \text{co-semisimplicial}\\ \text{pour } i\\ \text{variable}}}\ ) = \Ker(H^i(U,F) \rightrightarrows H^i(U^2,F)).
\]

\printedpage{20}
% head: A. Grothendieck : Dérivateurs, Ch. II, Édition des Universités de Montpellier II et Paris VII
% foot: Version 8 juillet 2001

Par hypothèse sur $U$, l'image de $\xi$ dans le troisième terme est nulle, donc $\xi$ appartient au premier terme. Mais pour $(*)$, le complexe $C(U^*,F)$ est le complexe \textit{trivial} des puissances de $I=\pi_0(U)$, à coefficients dans $M$ (si $F=M_{\mathcal{X}}$). Sa cohomologie \msp{page 32} en degrés $>0$ est nulle, donc $\xi=0$, q.e.d.

\noindent\textbf{Proposition 2.15.} (\footnote[39]{C'est déjà dans SGA 4, I 8.}) \textit{Soit $X$ un objet de} Cat\textit{. Conditions équivalentes :}
\begin{enumerate}[label={}, leftmargin=2em, itemsep=0pt]
\item \textit{a) $X$ est filtrant, i.e. $X$ est $0$-connexe et $X^o$ satisfait les conditions équivalentes du théorème 2.6 ($X^{o,\wedge}$ est localement totalement $0$-connexe).}
\item \textit{b) L'objet final de $X^\wedge$ est ind-représentable.}
\end{enumerate}
\noindent\textsc{Démonstration}. On sait qu'on a
\[
e_{X^\wedge} \simeq \varinjlim_{x \in X} x, \tag*{(2.15.1)}
\]
limite inductive prise dans $X^\wedge$. En effet, cela signifie aussi que pour tout objet $a$ dans $X$, on a
\[
\varinjlim_{x \in X} \ \underbrace{x(a)}_{\Hom_X(a,x)} \ \xrightarrow{\ \sim\ } e \qquad \text{(ensemble ponctuel)}
\]
Mais considérons l'objet $X' = a\backslash X \longrightarrow X$ coétale sur $X$, qui est donc Cat-cofibré sur $X$, on sait que l'on a
\[
\pi_0(X') \simeq \varinjlim_{x \in X} \pi_0(X'_x),
\]
où $X'_x = \Hom(a,x)$, donc on trouve
\[
\varinjlim_{x \in X} \Hom(a,x) = \pi_0(a\backslash X) = e
\]
car $a\backslash X$ est 0-connexe.

La formule (2.15.1) montre que si \msp{page 33} $X$ est filtrant, $e_{X^\wedge}$ est limite inductive filtrante de foncteurs représentables. On sait d'ailleurs qu'on peut réindexer par une catégorie \textit{ordonnée} filtrante (cf. SGA 4, I 8.1.6 (c'est dû à P. \textsc{Deligne}\,!)).

On a prouvé a $\Longrightarrow$ b, prouvons b $\Longrightarrow$ a. Supposons donc
\[
e_{X^\wedge} = \varinjlim_{i \in I} x_i,
\]
où $I$ est un (petit) ensemble ordonné filtrant, et $I \longrightarrow X$ un foncteur. Prouvons PS 1$'$ (forme droite), i.e. pour $a, b$ dans $X$, \texttt{[il]} existe $\vcenter{\hbox{\begin{tikzcd}[row sep=0.1em, column sep=1.2em]
a \arrow[dr] & \\
& x \\
b \arrow[ur] &
\end{tikzcd}}}$. En effet, comme $\varinjlim\limits_{I} \Hom(a,x_i) = e$, il s'ensuit que $\Hom(a,x_i) \neq \emptyset$ pour $i$ grand, de même $\Hom(b,x_i) \neq \emptyset$, donc il suffit de prendre $\vcenter{\hbox{\begin{tikzcd}[row sep=0.1em, column sep=1.2em]
a \arrow[dr] & \\
& x_i \\
b \arrow[ur] &
\end{tikzcd}}}$. Prouvons PS 2$_{\mathrm{d}}$, i.e. que pour \texttt{[une]} double flèche $a \overset{u}{\underset{v}{\rightrightarrows}} b$, \texttt{[il]} existe $b \longrightarrow x$ qui égalise.

\printedpage{21}
% head: A. Grothendieck : Dérivateurs, Ch. II, Édition des Universités de Montpellier II et Paris VII
% foot: Version 8 juillet 2001

\msp{page 34}
Soit $\xi_{i_0}$ dans $\Hom(b,x_{i_0})$, ses deux images dans $\Hom(a,x_{i_0})$ ne sont pas nécessairement égales, mais le deviennent dans la limite inductive, réduite à $e$, donc deviennent égales en composant avec $x_{i_0} \longrightarrow x_i$ pour $i$ assez grand, donc le composé $b \longrightarrow x_{i_0} \longrightarrow x_i$
\[
a \rightrightarrows b \longrightarrow x_{i_0} \longrightarrow x_i
\]
égalise la double flèche.

\noindent\textbf{Corollaire 2.16.} \textit{Soit $W$ un localiseur fondamental satisfaisant l'axiome des limites} Loc (8)\textit{. Si $X$ dans} Cat \textit{est une catégorie filtrante ou cofiltrante, alors $X$ est $W$-asphérique.}

Plus généralement et agréablement :

\noindent\textbf{Corollaire 2.17.} \textit{Soit $W$ comme dans 2.16. Alors toute catégorie filtrante est} totalement $W$-asphérique (\footnote[40]{cf. XII, §V pour cette notion, que j'ai beaucoup peiné à dégager et à comprendre\,!}), \textit{i.e. pour toute flèche $f : X \longrightarrow Y$ dans} Cat\textit{, l'objet $f_!(e_X)$ de $e_Y$ (défini aussi comme $y \longmapsto \pi_0(y\backslash X)$, et intervenant dans la factorisation canonique de $f$ en $ip$, avec $p$ étale et les cofibres} \texttt{[?]} \textit{de $p$ $0$-connexes)} \msp{page 35} \textit{est $W$-asphérique. En fait, $Y/f_!(e_{X^\wedge})$ est même une catégorie filtrante.}

En effet, le foncteur $f_!$ commute aux limites inductives, donc $f_! e_{X^\wedge}$ est limite inductive filtrante des $f_!(x_i) = f(x_i)$. Donc $Y/f_!(e_{X^\wedge})$ est limite inductive filtrante des $Y/f(x_i)$. Donc c'est même une catégorie filtrante, a fortiori elle est $W$-asphérique.

Donc finalement, on trouve des exemples non triviaux de catégories totalement $W$-asphériques, i.e. des exemples donc des catégories sans objet final. % [sic] donc -> dont (?), the conjunction "donc" makes no sense here

\noindent\textbf{Corollaire 2.18.} (\footnote[41]{déjà dans SGA 4, I 8.}) \textit{Soit $X$ dans} Cat\textit{, $F$ dans $X^\wedge$. Pour que $F$ soit ind-représentable il faut et il suffit que $X/F$ soit filtrante.}

En effet, par (2.15), $X/F$ est filtrante si et seulement si l'objet final $F$ de $X/F$ provient d'un ind-objet de $X/F$, mais cela signifie aussi que $F$ est la limite inductive dans $X^\wedge$ d'un ind-objet de $X$, i.e. que $F$ est ind-représentable.

\msp{page 36}
\noindent\textbf{Question.} \textit{Si $X$ dans} Cat \textit{est totalement $W_\infty$-asphérique, i.e. les $f_!(e_{X^\wedge})$ sont $W_\infty$-asphériques pour toute flèche $f : X \longrightarrow Y$ dans} Cat\textit{, ou encore pour toute telle flèche, avec les $y\backslash Y$ $0$-connexes ($y \in \Ob Y$), $Y$ est $W_\infty$-asphérique - cela implique-t-il que $X$ soit filtrante ?} Ça ne m'étonnerait pas - et même en supposant seulement les $f_!(e_{X^\wedge})$ 1-connexes (\footnote[42]{Ils le sont toujours (cf. XII, p. 101), c'est 2-\textit{connexe} qu'il faut lire ici\,!}) (ils sont 0-connexes, dès que $X$ l'est). Mais je ne veux pas m'attarder à présent sur cette question.

La réponse à la question précédente est négative, comme on voit par la :

\printedpage{22}
% head: A. Grothendieck : Dérivateurs, Ch. II, Édition des Universités de Montpellier II et Paris VII
% foot: Version 8 juillet 2001

\noindent\textbf{Proposition 2.19.} (\footnote[43]{\textbf{Faux.}}) \textit{Soit $X$ dans} Cat \textit{telle que $\mathit{diag}_X : X \longrightarrow X \times X$ soit cofinal} (\footnote[44]{cf. définition 3.3, p. 45.})\textit{, i.e. telle que $X^o$ soit totalement $0$-connexe. Alors $X$ est totalement $W_\infty$-asphérique (donc aussi totalement $W$-asphérique pour tout localiseur fondamental $W$ tel que $W \supseteq W_\infty$).}

Soit en effet un foncteur cofinal $f : X \longrightarrow Y$, il faut prouver que $Y$ est $W_\infty$-asphérique. Or on sait que totalement 0-connexe implique $W_\infty$-asphérique (\footnote[45]{C'est \textbf{faux.}}) (cf. proposition 2.3, p. 17-18), donc il suffit de voir que $Y^o$ est totalement 0-connexe, i.e. $Y \xrightarrow{d_Y} Y \times Y$ cofinal. Or on a le diagramme commutatif
\[
\begin{tikzcd}[row sep=large, column sep=large]
X \arrow[r, "d_X"] \arrow[d, "f"'] & X \times X \arrow[d, "f \times f"] \\
Y \arrow[r, "d_Y"] & Y \times Y,
\end{tikzcd}
\]
où par hypothèse $f$, donc aussi $f \times f$, et $d_X$ sont cofinaux. Donc $\varphi = (f \times f) \circ p_X$ \texttt{[plutôt} $\varphi = (f \times f) \circ d_X$\texttt{]} est cofinal, et comme $\varphi = d_Y f$ et $f$ est aussi cofinal, on conclut que $d_Y$ l'est, q.e.d.

\noindent\textbf{Corollaire.} $\mathbf{\Delta}$ \textit{étant totalement $0$-connexe, $\mathbf{\Delta}^o$ est totalement $W_\infty$-asphérique, pourtant elle n'est pas filtrante.}

\noindent\textbf{Question.} \textit{$X$ totalement $W_\infty$-asphérique $\Longrightarrow$ $X^o$ totalement $0$-connexe ? Ou du moins implique PS 1$'$ ?}

\msp{page 36bis}
\noindent\textbf{Contrexemple 2.20} à la proposition fausse 2.19.

Soit $X$ un objet de Cat, rajoutons-lui la catégorie finale $e$ comme fermé, i.e. considérons un objet $X'$ de Cat dans lequel $X$ soit un sous-objet ouvert, $e$ le fermé complémentaire :
\[
X \overset{i}{\hookrightarrow} X' \overset{j}{\hookleftarrow} e.
\]
On sait (cf. XII) qu'une telle opération est déterminée, à isomorphisme près, par la donnée d'un objet
\[
H \in \Ob X^\wedge
\]
en posant
\[
H(x) = \Hom_{X'}(x,e),
\]
la catégorie $X/e$ s'identifiant (à isomorphisme près) à $X/H$. Elle est donc 0-connexe si et seulement si l'objet $H$ de $X^\wedge$ est 0-connexe. Cela signifie aussi que le morphisme $i$ est cocofinal, i.e. que $i^o : X^o \longrightarrow X'^o$ est cofinal. Supposons maintenant $X$ $W_\infty$-asphérique et totalement 0-connexe, i.e. l'actuel Loc (6), cf. XVI. La proposition 2.19 affirme que, sous ces conditions, $X'$ (qui est lui aussi totalement 0-connexe) est même $W$-asphérique. Il s'ensuivrait donc que $i \in W_\infty$. Mais par XII, th. 4, cor. 4 (p. 149), il s'ensuivrait que
\printedpage{23}
% head: A. Grothendieck : Dérivateurs, Ch. II, Édition des Universités de Montpellier II et Paris VII
% foot: Version 8 juillet 2001
\msp{page 36ter}
$X/e$ est $W_\infty$-asphérique, i.e. $H$ $W_\infty$-asphérique. Ainsi, si $X$ est totalement 0-asphérique, tout objet 0-connexe de $X^\wedge$ serait $\infty$-connexe. Mais c'est manifestement faux, p.ex. si on prend pour $X$ une catégorie-test pour $W_\infty$, p.ex. $X = \mathbf{\Delta}$.

Cela me fait soupçonner que la conjecture
\[
X \text{ totalement } W_\infty\text{-asphérique} \Longrightarrow X \text{ filtrante}
\]
pourrait être valable.

\end{document}
